In order to evaluate the impact of structural aliasing on the Chronos-Bolt framework and its forecasting capabilities, we conducted a series of experiments using synthetic time-series data.

To observe the predicted values generated by Chronos-Bolt, we use three type of signal generators: a single-armonic sinusoidal generator, a modified TSMixup generator and a KernelSynth generator. The last two are based on the methods proposed used in Chronos training.\cite{ansariChronosLearningLanguage2024}

\subsubsection{Single-Armonic Sinusoidal Generator}
This generator produces a simple sinusoidal signal with a single frequency component. The generated signal can be expressed mathematically as:

\begin{equation}
    x(t) = A \cdot \sin(2\pi f t + \phi)
\end{equation}

The continuous signal is then sampled at a specified rate to create a discrete time-series. The parameters of the generator, such as amplitude, frequency, and phase, can be adjusted to create different sinusoidal signals for testing purposes.

\subsubsection{Light TSMixup Generator}
We propose a modified version of the TSMixup method descriped in Chronos synthetic data augmentation.\cite{ansariChronosLearningLanguage2024} The Light TSMixup generator creates synthetic time-series data by combining multiple sinusoidal signals weighted by sampling from a symmetric Dirichlet distribution as in the original TSMixup method.

Instead of using a dataset of different time-series signals, we simplified the pool by using only continuos sinusoidal waves with different frequencies, amplitude, and phase. The resulting continuous signal is then sampled at a specified rate to create a discrete time-series. 

The parameters of the generator, such as the maximum number of sinusoidal components, and the sampling frequency, can be adjusted to create different synthetic signals for testing purposes. The pseudo-code for the Light TSMixup generator is listed in \autoref{sec:appB}

\subsubsection{KernelSynth Generator}
